% Generated draft: requires mathematical and source review.
\documentclass[11pt]{article}
\usepackage[margin=25mm]{geometry}
\usepackage{fontspec}
\setmainfont{Latin Modern Roman}
\usepackage{amsmath,amssymb,mathtools,unicode-math}
\setmathfont{Latin Modern Math}
\usepackage{enumitem,needspace,xurl,hyperref}
\setlength{\parindent}{0pt}
\setlength{\parskip}{5pt}
\setlength{\emergencystretch}{3em}
\newcommand{\N}{\mathbb{N}}
\newcommand{\Z}{\mathbb{Z}}
\newcommand{\Q}{\mathbb{Q}}
\newcommand{\R}{\mathbb{R}}
\newcommand{\C}{\mathbb{C}}
\newcommand{\F}{\mathbb{F}}
\newcommand{\A}{\mathbb{A}}
\newcommand{\PP}{\mathbb{P}}
\newcommand{\E}{\mathbb{E}}
\newcommand{\eps}{\varepsilon}
\newcommand{\dd}{\,\mathrm{d}}
\begin{document}
\section*{UnsolvedMath ID: 30002028}
\textbf{Source problem code:} catalogue entry
\subsection{Definitions and mathematical statement}
Let \(H=(V,E)\) be a finite simple \(r\)-uniform hypergraph: \(E\subseteq\binom{V}{r}\). Interpret the source's parameter \(n\) as \(|V|\). The complete \(r\)-uniform hypergraph \(K^r_\ell\) has \(\ell\) vertices and every \(r\)-element subset as an edge. Under the unlabelled, non-induced counting convention, define
\[
N_\ell(H)=\bigl|\{S\in\tbinom{V}{\ell}:\tbinom{S}{r}\subseteq E\}\bigr|.
\]
The source does not specify its counting convention; this convention is an explicit interpretation.

A complete \(\ell\)-partite \(r\)-uniform hypergraph with equal part size \(t\) has pairwise disjoint vertex classes \(V_1,\ldots,V_\ell\), each of size \(t\), and as edges all \(r\)-element sets meeting each class in at most one vertex. Containment here means subhypergraph containment: additional edges of \(H\) on these vertices are permitted.

For the meaningful exponent range \(r\geq2\), the supplied question, with the above conventions, is whether
\[
\begin{gathered}
\forall\ell\in\Z_{\geq1}\ \forall r\in\Z_{\geq2}\ \forall c\in\R_{>0}\
\exists\gamma=\gamma(\ell,r,c)>0\ \forall n\in\Z_{\geq1}\
\forall H=(V,E),\quad |V|=n,\ E\subseteq\tbinom{V}{r}:\\
N_\ell(H)\geq c n^\ell\ \Longrightarrow\
\exists t\in\Z_{\geq1}\ \exists V_1,\ldots,V_\ell\subseteq V,\\
V_i\cap V_j=\varnothing\ (i\ne j),\qquad |V_i|=t\ (1\leq i\leq\ell),\\
t\geq\gamma(\log n)^{1/(r-1)},\qquad
\bigl\{A\in\tbinom{\bigcup_{i=1}^{\ell}V_i}{r}:\ |A\cap V_i|\leq1\text{ for all }i\bigr\}\subseteq E.
\end{gathered}
\]
Here \(\log\) denotes the natural logarithm; changing its fixed base only changes \(\gamma\).

The supplied source actually quantifies over all positive integers \(r\), including \(r=1\), where \(1/(r-1)\) is undefined. Thus the display records only the portion with a defined exponent, not a resolution of the full supplied wording. No hypothesis \(\ell\geq r\) is supplied; when \(\ell<r\), both complete hypergraphs in these definitions have no edges. The universal quantifier over \(n\) above is a literal interpretation of the source's implicit order parameter; the source does not explicitly distinguish all \(n\) from sufficiently large \(n\). These qualifications require review.
\subsection{Short English statement}
The question asks whether a positive density of complete \(r\)-uniform hypergraphs on \(\ell\) vertices forces a complete \(\ell\)-partite subhypergraph whose equal parts have size at least a constant times \((\log n)^{1/(r-1)}\). The constant may depend on \(\ell,r,c\), but not on \(n\) or \(H\). The supplied wording includes an undefined case \(r=1\) and leaves several conventions implicit. This is a statement-only draft requiring human mathematical review; the supplied open label is not independently verified.
\subsection{Sources}
\begin{itemize}
\item Local UnsolvedMath registry, permanent ID \(30002028\), title ``Supersaturation Blow-Ups in Uniform Hypergraphs,'' Oberwolfach Reports Open Problems collection; exact source problem code \texttt{catalogue entry}. No individual proposer is supplied. The record gives proposed year 2012.
\item Supplied source citation: \emph{Mini-Workshop: Hypergraph Tur\'an Problem (2013)}, Oberwolfach Reports, DOI \texttt{10.4171/owr/2012/19}, pp.\ 1149--1182. \url{https://doi.org/10.4171/owr/2012/19}. The supplied background describes these reports as documenting mathematical workshops and their open-problem discussions.
\item Supplied literature source, ``Oberwolfach report'': \url{https://ems.press/content/serial-article-files/46390?nt=1}.
\item Supplied literature source, ``Finding partite hypergraphs efficiently'': \url{https://www.sciencedirect.com/science/article/pii/S0020019026000050}.
\item Literature assessment dated exactly \texttt{2026-08-21}: the supplied curation labels the problem open and says that the stronger conjectured logarithmic blowup exponent from a positive density of complete partite copies remains unresolved, while recent algorithmic work reaches the Erd\H{o}s-order exponent but not this copy-count conjecture. This is supplied dated curation, not an independently verified assessment. Its own qualification requests independent verification of time-sensitive status and cited claims.
\item The registry labels the statement \texttt{corrected\_verified} and supplies a claim of manual semantic, typographic, and LaTeX review of a meaning-preserving cleanup. These are supplied provenance labels, not mathematical verification by this draft.
\end{itemize}
\end{document}
